\documentclass[12pt,addpoints]{evalua}
\grado{2$^\circ$ de Secundaria}
\cicloescolar{2026-2027}
\materia{Matemáticas 2}
\unidad{1}
\title{Examen Final}
\aprendizajes{
      \item Resuelve problemas de multiplicación y división con números enteros, fracciones y decimales positivos y negativos.
      \item Resuelve problemas de potencias con exponente entero y aproxima raíces cuadradas.
      \item Resuelve problemas que impliquen el uso de la notación científica.
      \item Calcula porcentajes de cantidades.
}
\author{Prof.: Julio César Melchor Pinto}
\begin{document}%
\begin{multicols}{2}
	\tableofcontents
\end{multicols}\newpage
\begin{questions}
	\section{Unidad 1}
	\subsection   {Cálculos numéricos}
	\subsubsection{Suma de números}
	\question{\include*{../questions/ques_calc_num-suma_numeros_1}}
	\subsubsection{Resta de números}
	\question{\include*{../questions/ques_calc_num-resta_numeros_1}}
	\subsubsection{Multiplicación de números}
	\question{\include*{../questions/ques_calc_num-multiplicacion_numeros_1}}
	\subsubsection{División de números}
	\question{\include*{../questions/ques_calc_num-division_numeros_1}}
	\subsubsection{Resolución de Problemas}
	\question{\include*{../questions/ques_calc_num-resolucion_problemas_1}}
	\subsection   {Fracciones}
	\subsubsection{Clasificación de fracciones}
	\question{\include*{../questions/ques_fracciones-clasificacion_1}}
	\subsubsection{Representación de fracciones}
	\question{\include*{../questions/ques_fracciones-representacion_1}}
	\subsubsection{Nombre de fracciones}
	\question{\include*{../questions/ques_fracciones-nombre_1}}
	\subsubsection{Fracciones en la recta numérica}
	\question{\include*{../questions/ques_fracciones-recta_numerica_1}}
	\subsubsection{Conversión de fracciones}
	\question{\include*{../questions/ques_fracciones-conversion_1}}
	\subsection   {Fracciones MCD y MCM}
	\subsubsection{Comparación de fracciones}
	\question{\include*{../questions/ques_fracciones-comparacion_1}}
	\subsubsection{Fracciones equivalentes}
	\question{\include*{../questions/ques_fracciones-equivalentes_1}}
	\subsubsection{M.C.D y M.C.M}
	\question{\include*{../questions/ques_fracciones-mcd_mcm_1}}
	\subsubsection{Simplificación de fracciones}
	\question{\include*{../questions/ques_fracciones-simplificacion_1}}
	\subsubsection{Resolución de problemas}
	\question{\include*{../questions/ques_fracciones-resolucion_problemas_1}}
	\subsection   {Números decimales}
	\subsubsection{Ubicación en la recta numérica}
	\question{\include*{../questions/ques_decimales-recta_numerica_1}}
	\subsubsection{Porcentajes a decimal}
	\question{\include*{../questions/ques__decimales-porcentajes-decimal_1}}
	\subsubsection{Operaciones con múltiplos de 10}
	\question{\include*{../questions/ques_decimales-operaciones-multiplos_10_1}}
	\subsubsection{Conversión de fracciones a decimales}
	\question{\include*{../questions/ques_decimales-conversion_fracc_decimal_1}}
	\subsubsection{Conversión de decimales a fracciones}
	\question{\include*{../questions/ques_decimales-conversion_decimal_fracc_1}}
	\subsection   {Conceptos básicos de números negativos}
	\subsubsection{Ubicación en la recta numérica}
	\question{\include*{../questions/ques_negativos-recta_numerica_1}}
	\subsubsection{Comparación de negativos}
	\question{\include*{../questions/ques_negativos-comparacion_1}}
	\subsubsection{Oraciones con negativos}
	\question{\include*{../questions/ques_negativos-oraciones_1}}
	\subsubsection{Determina el signo en suma y resta con negativos}
	\question{\include*{../questions/ques_negativos-signo_suma_resta_1}}
	\subsubsection{Determina el signo multiplicación y división con negativos}
	\question{\include*{../questions/ques_negativos-signo_multiplicacion_division_1}}
	\subsection   {Operaciones con números negativos}
	\subsubsection{Comparación de negativos}
	\question{\include*{../questions/ques_negativos-comparacion_2}}
	\subsubsection{Suma y resta con negativos}
	\question{\include*{../questions/ques_negativos-suma_resta_1}}
	\subsubsection{Multiplicación y división con negativos}
	\question{\include*{../questions/ques_negativos-multiplicacion_division_1}}
	\subsubsection{Potencias con numeros negativos}
	\question{\include*{../questions/ques_negativos-potencias_1}}
	\subsubsection{Jerarquía de operacione}
	\question{\include*{../questions/ques_negativos-jerarquia_1}}

	\section{Unidad 2}
	\subsection   {Círculo}
	\subsubsection{Resolución de problemas}
	\question{\include*{../questions/}}
	\subsubsection{Radio, Diámetro, Perímetro y Área de un círculo}
	\question{\include*{../questions/}}
	\subsection   {Polígonos y circunferencias}
	\subsubsection{Ángulos interiores}
	\question{\include*{../questions/}}
	\subsubsection{Ángulos centrales y exteriores}
	\question{\include*{../questions/}}
	\subsubsection{Ángulos centrales e inscritos}
	\question{\include*{../questions/}}
	\subsubsection{Arco de una circunferencia}
	\question{\include*{../questions/}}
	\subsubsection{Área de un sector circular}
	\question{\include*{../questions/}}
	\subsection   {Figuras y cuerpos geométricos}
	\subsubsection{Perímetro y Área}
	\question{\include*{../questions/}}
	\subsubsection{Resolución de problemas}
	\question{\include*{../questions/}}
	\subsubsection{Área lateral, Área total y Volumen}
	\question{\include*{../questions/}}
	\subsection   {Monomios y polinomios}
	\subsubsection{Lenguaje algebraico}
	\question{\include*{../questions/}}
	\subsubsection{Suma de monomios y polinomios}
	\question{\include*{../questions/}}
	\subsubsection{Resta de monomios y polinomios}
	\question{\include*{../questions/}}
	\subsubsection{Operaciones combinadas}
	\question{\include*{../questions/}}
	\subsubsection{Perímetro de figuras geométricas}
	\question{\include*{../questions/}}
	\subsection   {Operaciones con monomios y polinomios}
	\subsubsection{Suma, resta y multiplicación de exponentes}
	\question{\include*{../questions/}}
	\subsubsection{Multiplicación y división de monomios y polinomios}
	\question{\include*{../questions/}}
	\subsubsection{Áreas de figuras geométricas}
	\question{\include*{../questions/}}
	\subsection   {Sistema de unidades}
	\subsubsection{Unidades de longitud}
	\question{\include*{../questions/}}
	\subsubsection{Unidades de masa}
	\question{\include*{../questions/}}
	\subsubsection{Unidades de capacidad}
	\question{\include*{../questions/}}
	\subsubsection{Unidades de área y volumen}
	\question{\include*{../questions/}}

	\section{Unidad 3}
	\subsection   {Probabilidad y estadística}
	\subsubsection{Mediana y moda}
	\question{\include*{../questions/}}
	\subsubsection{Promedio}
	\question{\include*{../questions/}}
	\subsubsection{Interpretación de gráficas}
	\question{\include*{../questions/}}
	\subsubsection{Eventos mutuamente excluyentes}
	\question{\include*{../questions/}}
	\subsubsection{Eventos dependientes e independientes}
	\question{\include*{../questions/}}
	\subsection   {Razones y proporciones}
	\subsubsection{Relaciones proporcionales}
	\question{\include*{../questions/}}
	\subsubsection{Constante de proporcionalidad}
	\question{\include*{../questions/}}
	\subsubsection{Regla de correspondencia (ecuación)}
	\question{\include*{../questions/}}
	\subsubsection{Proporción directa e inversa}
	\question{\include*{../questions/}}
	\subsubsection{Proporciones compuestas}
	\question{\include*{../questions/}}
	\subsection   {Sucesiones aritméticas}
	\subsubsection{Completando la sucesión}
	\question{\include*{../questions/}}
	\subsubsection{Diferencia de una sucesión}
	\question{\include*{../questions/}}
	\subsubsection{Término enésimo}
	\question{\include*{../questions/}}
	\subsubsection{Término general}
	\question{\include*{../questions/}}
	\subsubsection{Término enésimo 2}
	\question{\include*{../questions/}}
	\subsection   {Ecuaciones lineales}
	\subsubsection{Lenguaje algebraico}
	\question{\include*{../questions/}}
	\subsubsection{Sustitución de valores}
	\question{\include*{../questions/}}
	\subsubsection{Ecuaciones de primer grado}
	\question{\include*{../questions/}}
	\subsubsection{Resolución de problemas}
	\question{\include*{../questions/}}
	\subsection   {Sistemas de ecuaciones}
	\subsubsection{Método de eliminación}
	\question{\include*{../questions/}}
	\subsubsection{Método de sustitución}
	\question{\include*{../questions/}}
	\subsubsection{Método de igualación}
	\question{\include*{../questions/}}
	\subsubsection{Sistema de ecuaciones}
	\question{\include*{../questions/}}
\end{questions}
\end{document}